\documentclass[11pt]{article}

% ============================================================
% Packages
% ============================================================
\usepackage[utf8]{inputenc}
\usepackage[margin=1in]{geometry}
\usepackage{amsmath,amssymb,amsthm}
\usepackage{graphicx}
\usepackage{booktabs}
\usepackage{xcolor}
\usepackage{hyperref}
\usepackage{tikz}
\usepackage{enumitem}
\usepackage{pgfplots}

\pgfplotsset{compat=1.18}

% ============================================================
% Document Formatting
% ============================================================
\setlength{\parskip}{0.6em}
\setlength{\parindent}{0pt}

% ============================================================
% List Formatting
% ============================================================
\setlist[enumerate]{
    leftmargin=2em,
    itemsep=0.25em
}

% ============================================================
% Solution Environment
% ============================================================
\newenvironment{solution}
    {\par\noindent\textbf{Solution.}\par}
    {\par}

% ============================================================
% Link Styling
% ============================================================
\hypersetup{
    colorlinks=true,
    linkcolor=blue,
    urlcolor=cyan
}

% ============================================================
% Metadata
% ============================================================
\title{Ordinary Differential Equations: Abridged Notes}
\author{Chris Ritter}
\date{September 29, 2026}

% ============================================================
% Document
% ============================================================

\begin{document}
\maketitle
\begin{center}
\small\emph{These abridged notes were reconstructed afterward because I did not take
direct notes during class. They are not a record of the actual lecture.}
\end{center}
\tableofcontents
\newpage

\section{Topic Summary}
Complex roots of a real constant-coefficient characteristic equation
still give real solutions. A conjugate pair \(r=\alpha\pm i\beta\)
produces \(e^{\alpha t}\cos\beta t\) and
\(e^{\alpha t}\sin\beta t\). The real part \(\alpha\) controls the
amplitude: negative gives decay, positive gives growth, and zero
gives constant-amplitude oscillation. Initial conditions fix the two
constants.

\section{From Complex Roots to Real Solutions}
For \(ay''+by'+cy=0\), solve \(ar^2+br+c=0\). If the roots are
\(\alpha\pm i\beta\), with \(\beta>0\), Euler's formula gives the
real general solution
\[
\boxed{y(t)=e^{\alpha t}
\bigl(C_1\cos\beta t+C_2\sin\beta t\bigr).}
\]
The oscillation period is \(2\pi/\beta\). The factor
\(e^{\alpha t}\) is an amplitude envelope, so the sign of
\(\alpha\) determines long-term growth or decay.

\subsection{Example: Damped Oscillation}
\textbf{Adapted from Boyce et al., 11th ed., \S3.3, Problem 8.}
Solve \(y''+6y'+13y=0\).
\begin{solution}
The characteristic equation has roots
\[
r=\frac{-6\pm\sqrt{36-52}}{2}=-3\pm2i.
\]
Therefore
\[
\boxed{y(t)=e^{-3t}\bigl(C_1\cos2t+C_2\sin2t\bigr).}
\]
The solution oscillates with period \(\pi\), while its amplitude
shrinks like \(e^{-3t}\).
\end{solution}

\subsection{Example: Growing Oscillation with Initial Data}
\textbf{Adapted from Boyce et al., 11th ed., \S3.3, Problem 13.}
Solve
\[
y''-2y'+5y=0,\qquad y(\pi/2)=0,\qquad y'(\pi/2)=2.
\]
\begin{solution}
The roots are \(1\pm2i\). Center the solution at
\(t_0=\pi/2\) by writing \(s=t-\pi/2\):
\[
y=e^s(A\cos2s+B\sin2s).
\]
At \(s=0\), \(y=A=0\) and \(y'=A+2B=2\), hence \(B=1\). Thus
\[
\boxed{y(t)=e^{t-\pi/2}\sin\bigl(2(t-\pi/2)\bigr).}
\]
It crosses zero periodically, but its oscillation envelope grows as
\(t\) increases because the characteristic roots have positive real
part.
\end{solution}
\end{document}
